\section{Walls in general}\label{sec:walls}

In this section $F$ is a reversible cellular automaton of radius $1$ on a finite alphabet $\A$. A nonempty word $W$ is \emph{fixed in every context} if $F(c)$ contains $W$ at every position at which $c$ contains $W$.

\begin{lemma}\label{lem:exact}
Under the hypothesis of Theorem~\ref{thm:walls}, $F$ neither creates nor destroys occurrences of $W$: $F(c)$ contains $W$ at position $i$ if and only if $c$ does.
\end{lemma}

\begin{proof}
Fix $n$ and let $P_n$ be the finite set of configurations of spatial period $n$. $F$ maps $P_n$ bijectively to itself, being injective and commuting with $\sh^n$. Write $N(c)$ for the number of positions $i\in\{0,\dots,n-1\}$ at which $c$ contains $W$. By hypothesis $N(F(c))\ge N(c)$ for every $c\in P_n$, and $\sum_{c\in P_n}N(F(c))=\sum_{c\in P_n}N(c)$ because $F$ permutes $P_n$. Hence $N(F(c))=N(c)$, so on $P_n$ the occurrences of $W$ in $F(c)$ are exactly those of $c$. If some configuration $c$ had an occurrence of $W$ in $F(c)$ at a position $i$ where $c$ has none, the same would hold for the periodic configuration obtained by repeating a long enough window of $c$ around $i$, since both properties depend only on $c_{i-1},\dots,c_{i+|W|}$.
\end{proof}

\begin{proof}[Proof of Theorem~\ref{thm:walls}]
Let $k=|W|$. A \emph{car} of $c$ is a position $i$ with $c$ containing $W$ at $i-k$ and at $i+1$ and $c_i\in C$; by the definition of $C$ no occurrence of $W$ covers $i$. By Lemma~\ref{lem:exact} the occurrences of $W$ in $F(c)$ are those of $c$, and the cells of each occurrence keep their values, so the neighbours of a car keep their values and the car evolves by a map $\varphi\colon C\to\A$, $\varphi(s)=F(W\,s\,W)$ at the middle position, independent of the rest of $c$. The image $\varphi(s)$ lies in $C$, because an occurrence of $W$ in $W\,\varphi(s)\,W$ other than at its ends would be an occurrence created by $F$. And $\varphi$ is injective, because $F$ is injective on the periodic configurations $(W\,s)^{\infty}$, which differ only at their cars. So $\varphi$ is a permutation of $C$, and $F$ maps cars to cars. Conversely every car of $F(c)$ is a car of $c$: its position is flanked by the same occurrences of $W$, and if the symbol of $c$ there were outside $C$, then $c$ would contain an occurrence of $W$ covering it, which $F$ would preserve. So cars are invariant, and $F$ acts on every car by the same $\varphi$; this uses that the neighbours of a car lie in occurrences of $W$, whose images under $F$ do not depend on the context, which is where the hypothesis ``fixed in every context'' enters. A \emph{train with $m$ cars} is a chain of cars $i_1<\dots<i_m$ with $i_{j+1}=i_j+k+1$, such that $c$ contains $W$ at $i_1-2k$ and at $i_m+k+1$ (two copies of $W$ with no car between them at each end), and that is maximal with these properties. Trains are determined by the occurrences of $W$ and the cars, so they are invariant under $F$. As in the proof of Theorem~\ref{thm:finite}, permuting the cars of every $m$-car train by $\sigma\in S_m$ defines a block map $\Psi_m(\sigma)$ of radius at most $m(k+1)+2k+1$, which is an automorphism commuting with $F$ because $F$ applies the same $\varphi$ to every car. With two distinct symbols $a,b\in C$, the periodic configuration whose period is one $m$-car train with cars $(b,\dots,b,a,b,\dots,b)$ shows that $\sigma$ is recovered from $\Psi_m(\sigma)$. The map $\Psi_m(\sigma)$ preserves the occurrences of $W$, since any occurrence covering a car lies inside $W\,s\,W$ and the new symbol is again in $C$; hence it preserves cars and trains. Maps for different $m$ act on disjoint cells and commute, and the resulting homomorphism of $\bigoplus_mS_m$ is injective, as one sees by evaluating on periodic configurations $(W\,W\,s_1\,W\,s_2\cdots W\,s_m)^{\infty}$ whose only trains have $m$ cars.
\end{proof}

The train definition here, bounded by two adjacent copies of $W$, differs from the one used for rule $54$ in Theorem~\ref{thm:finite}, where outer walls of length at least $3$ suffice; both definitions work, and the general one is what is needed for words that overlap themselves.

\begin{table}[t]
\caption{Second-order elementary automata with a word $W$ fixed in every context, one rule per conjugacy class, with the shortest such $W$ (cells written $0$--$3$). The words are the fixed walls (i)--(viii) of Takesue \cite[Fig.~2, Table~III]{Takesue1989}, and the $31$ classes are those of the rules listed there for these walls. Classes are taken under reflection and the cellwise complement $\kappa(x,y)=(1-x,1-y)$, which conjugates $F_f$ to $F_{f'}$ with $f'(a,b,c)=f(1-a,1-b,1-c)$, as in \cite[Table~I]{Takesue1989}; Wolfram's complement $1\oplus f(1-a,1-b,1-c)$ is not a conjugacy of second-order automata \cite[p.~377]{Takesue1989}, and for instance $110$R and $137$R have different numbers of periodic points. The property is constant on each class. For each listed rule $|C|\ge2$. Of the $57$ classes in the last row, only $133=\{133,161\}$ has a fixed word of greater length ($W=0330$, $k=4$, the wall (ix) of \cite{Takesue1989}; Section~\ref{cl:sec:133}).}
\label{tab:walls}
\centering\small
\begin{tabular}{l}
\toprule
$W=0$: 0, 4, 12 \qquad $W=3$: 1, 2, 3, 18, 19, 33, 35, 50, 51\\
$W=30$: 9, 10, 11, 131, 139 \qquad $W=03$: 129 \qquad $W=030$: 73 \qquad $W=303$: 147 \qquad $W=00$: 36\\
$W=33$: 5, 6, 7, 22, 23, 37, 38, 39, 54, 55\\
\midrule
no fixed word of length $\le3$ ($57$ classes): 13--15, 24--31, 41, 43, 45--47, 57--63, 75, 77, 79, 90, 91,\\
\quad 94, 95, 105, 107, 109, 111, 123, 126, 127, 133, 135, 141, 143, 151, 153, 155, 157, 159, 165, 167,\\
\quad 175, 179, 183, 187, 189, 191, 219, 223, 255\\
\bottomrule
\end{tabular}
\end{table}

\begin{table}[t]
\caption{Spread of a single changed cell after $300$ steps, mean over $40$ random configurations of length $801$, and number of trials in which the difference disappeared.}
\label{tab:spread}
\centering
\begin{tabular}{lrrr}
\toprule
rule & leftward spread & rightward spread & healed\\
\midrule
54R & 15 & 21 & 0\\
110R & 195 (min 167) & 41 (min 22) & 0\\
30R & 98 (min 72) & 300 (min 299) & 0\\
\bottomrule
\end{tabular}
\end{table}


\begin{proof}[Proof of Corollary~\ref{cor:48}]
For each rule of Table~\ref{tab:walls} the listed word $W$ is fixed in every context, and the set $C$ of Theorem~\ref{thm:walls} has at least two elements; both are finite checks of the local rule over the $4^2$ contexts of $W$ and the four symbols $s$ in $W\,s\,W$. Theorem~\ref{thm:walls} applies, and by the conjugacy invariance of Section~\ref{sec:prelim} it applies to every rule of the $31$ classes, $81$ rules in all. The class of $133$ is Proposition~\ref{cl:prop:133}(1), and the last statement is Proposition~\ref{cl:prop:133}(3).
\end{proof}

The walls also confine the spreading of differences: in random configurations of $54$R a single changed cell produces a difference that spreads only a few cells, while in $110$R and $30$R it spreads in both directions at linear speed (Table~\ref{tab:spread}).

\subsection{The class of 133: a wall of length 4}\label{cl:sec:133}

Table~\ref{tab:walls} lists the fixed words of length at most $3$. One more class has a fixed word, of length $4$. It is the wall (ix) of Takesue, who notes that it is not detected with fewer than four cells \cite[Table~III and p.~391]{Takesue1989}. In the next proposition the fixed word of part~(1) is his; the automorphisms of parts~(1) and~(2) and the completeness statement~(3) are what is added. The examination in \cite[Section~4]{Takesue1989} covers the words of at most $6$ cells, and it is accompanied by the conjecture, stated there without proof, that $4$ cells suffice to detect a local conservation law of a second-order elementary automaton. Part~(3) concerns fixed words only, and for them it gives a proof for every length.

\begin{proposition}[the fixed word of $133$R: {\cite[p.~391]{Takesue1989}}]\label{cl:prop:133}
Let $f\in\{133,161\}$ and $F=F_f$.
\begin{enumerate}
\item The word $W=0330$ is fixed by $F_{133}$ in every context, and $W=3003$ is fixed by $F_{161}$ in every context. In both cases the set $C$ of Theorem~\ref{thm:walls} is all of $\A$, so $\bigoplus_mS_m$ embeds in $\Aut(X_F)$.
\item For $f=133$, let $D$ be the block map of radius $4$ that replaces the middle cell $s$ of every occurrence of $W\,s\,W$ by $F(W\,s\,W)$ at that cell and fixes every other cell. Then $D\in\Aut(X_F)$, $D$ acts on each such cell as the $4$-cycle $0\mapsto1\mapsto3\mapsto2\mapsto0$, $D^4=\mathrm{id}\ne D^2$, and $D\notin\langle\sh,F\rangle$.
\item Every word fixed by some $F_{f'}$ in every context lies in $\{0,3\}^k$. Among the $57$ classes with no such word of length at most $3$, the class of $133$ is the only one with such a word of any length.
\end{enumerate}
\end{proposition}

\begin{proof}
(1) A cell of a word fixed in every context keeps its value, so by \eqref{eq:F} its two bits agree and $f(y_{j-1},y_j,y_{j+1})=0$ at each of its positions $j$, the outer neighbour bits being arbitrary. For $133$, $f(a,b,c)=1$ exactly on $000$, $010$ and $111$. The $y$-row of $W=0330$ is $0110$, and the triples $(u,0,1)$, $(0,1,1)$, $(1,1,0)$ and $(1,0,v)$, $u,v\in\{0,1\}$, all avoid $000$, $010$ and $111$; so $W$ is fixed in all $16$ contexts. Since $161=\kappa(133)$ and $\kappa$ maps $0330$ to $3003$, the statement for $161$ follows. Finally $W\,s\,W$ contains $W$ only at its two ends for each $s\in\A$, so $C=\A$. Theorem~\ref{thm:walls} applies.

(2) The neighbours of the middle cell $s$ of $W\,s\,W$ have $y$-bit $0$, so it evolves by $(x,y)\mapsto(y,x\oplus f(0,y,0))=(y,x\oplus1)$, the $4$-cycle $0\to1\to3\to2\to0$. So $D$ is the map that applies the car permutation $\varphi$ of the proof of Theorem~\ref{thm:walls} to every car. As there, every length-$4$ window covering a car lies in $W\,s\,W$, and $W\,\varphi(s)\,W$ contains $W$ only at its ends, so $D$ preserves the occurrences of $W$ and the cars, $D^{-1}=D^3$, and $D$ commutes with $F$ because $F$ applies $\varphi$ to every car and the neighbours of a car lie in occurrences of $W$, whose images do not depend on the context. Since $\varphi$ is a $4$-cycle, $D^4=\mathrm{id}\ne D^2$. Suppose $D=\sh^aF^b$. The configuration $0^{\Z}$ contains no $W$, so $D$ fixes it; $F$ acts on constant configurations as the $4$-cycle $0\to1\to3\to2\to0$, and $\sh$ fixes them, so $4\mid b$. On $c=(0330\,0)^{\infty}$, every fifth cell is a car, so $D(c)=F(c)$, and therefore $\sh^aF^{b-1}c=c$. But $F^{b-1}c=(0330\,s)^{\infty}$ with $s=\varphi^{b-1}(0)\ne0$, as $b-1\not\equiv0\pmod4$, and no shift of it equals $c$. So $D\notin\langle\sh,F\rangle$.

(3) The first statement is the first sentence of the proof of (1). With $w$ the $y$-row, a fixed word of length $k\ge2$ is a word $w\in\{0,1\}^k$ with $f(u,w_0,w_1)=0$ and $f(w_{k-2},w_{k-1},v)=0$ for $u,v\in\{0,1\}$ and $f=0$ on its interior triples: a walk through $k-1$ states in a graph on the $4$ pairs $(w_j,w_{j+1})$, from a set of allowed initial pairs to a set of allowed final pairs. If such a walk exists, a shortest one visits each state at most once, so some fixed word has length at most $5$. A search of all lengths $k\le40$ for all $256$ rules therefore decides the question for every length; it finds fixed words for exactly the $81$ rules of Table~\ref{tab:walls} at length at most $3$, and, among the other rules, only for $133$ and $161$, at the lengths $4$, $7$, $8$ and every $k$ with $10\le k\le40$. A brute-force search over all of $\A^k$, $k\le4$, agrees with the reduction to $\{0,3\}^k$. The search is a finite computation, described in Appendix~\ref{app:proofs}.
\end{proof}

